\NSPage{007}%
\begin{EESegment}{EE-TGT-P007-001}
The estimates have to control every Cartesian space--time derivative of this residual.
\end{EESegment}

\begin{EESegment}{EE-TGT-P007-002}
First, the proof builds the axisymmetric background
\NSi{NS-F-P007-001}{(u_B,p_B)} from Proposition~\ref{prop:5.5}. Its residual in
\eqref{eq:5.41} is the negative divergence
of a stress supported in an annulus, apart from a remainder that vanishes to every order at the
singular point. The oscillating velocities produce the leading part of this stress when their
quadratic products are averaged, as Proposition~\ref{prop:7.5} shows. Corrections added one after
another make the residual left over from that step decay faster, as Proposition~\ref{prop:9.6}
shows. The fields are then added together and confined to a fixed region without losing
incompressibility or the leading growth of the velocity. This lets the residual extend to the smooth,
compactly supported force in Theorem~\ref{thm:1.1}, by Propositions~\ref{prop:10.1}
and~\ref{prop:9.9} and Lemma~\ref{lem:10.3}. Every estimate used below is proved for the viscous
equations. Figures~\ref{fig:5} and~\ref{fig:6} show how the whole construction fits together.
\end{EESegment}

\begin{EESegment}{EE-TGT-P007-003}
For every \NSi{NS-F-P007-002}{\nu>0}, any solution with viscosity one gives a solution with
viscosity \NSi{NS-F-P007-003}{\nu} after the following change of scale.
\NSFormula{NS-F-P007-004}{display}{none}
\[
 u_\nu(x,t)=\sqrt{\nu}\,u(x/\sqrt{\nu},t),\qquad
 p_\nu(x,t)=\nu p(x/\sqrt{\nu},t),\qquad
 f_\nu(x,t)=\sqrt{\nu}\,f(x/\sqrt{\nu},t).
\]
This change does not move the singular time. Equations~\eqref{eq:10.22}--\eqref{eq:10.23}
check both the equation and the way the energy changes under this scaling. The rest of this
outline uses \NSi{NS-F-P007-005}{\nu=1}.
\end{EESegment}

\begin{EESegment}{EE-TGT-P007-004}
\subsection{The concentrating leading field.}\label{sec:3.1}
The fixed velocity profiles \NSi{NS-F-P007-006}{E,U} and pressure profile
\NSi{NS-F-P007-007}{\Pi} from \eqref{eq:4.3} put the geometry of Section~\ref{sec:2.1} into the
leading field. The similarity coordinates \NSi{NS-F-P007-008}{(X,\eta)} from \eqref{eq:3.2}
keep track of the different rates at which the field contracts in the radial and axial directions.
They also write the growth of the velocity as powers of one concentration scale. This subsection
sets up that representation. The next one builds profiles whose remaining momentum residual can
be supplied by the pulses.
\end{EESegment}

\begin{EESegment}{EE-TGT-P007-005}
Set
\NSFormula{NS-F-P007-009}{display}{none}
\[
 \tau=1-t,\qquad A=\frac12+h,\qquad D=\frac12-h,
 \qquad 0<h<1/100,
\]
and keep \NSi{NS-F-P007-010}{h} fixed while the profiles are built.
\end{EESegment}

\begin{EESegment}{EE-TGT-P007-006}
In cylindrical coordinates,
\NSFormula{NS-F-P007-011}{display}{none}
\[
 x=(r\cos\theta,r\sin\theta,z),\qquad
 e_r=(\cos\theta,\sin\theta,0),\qquad
 e_\theta=(-\sin\theta,\cos\theta,0),\qquad
 e_z=(0,0,1).
\]
The leading field
\NSi{NS-F-P007-012}{u^{(0)}=u_r^{(0)}e_r+u_\theta^{(0)}e_\theta+u_z^{(0)}e_z}
is axisymmetric, meaning that its cylindrical components do not depend on
\NSi{NS-F-P007-013}{\theta}. Its two tangential components are the azimuthal component and the
axial component. Both point along directions tangent to the cylinders on which
\NSi{NS-F-P007-014}{r} is constant.
\end{EESegment}

\begin{EESegment}{EE-TGT-P007-007}
Use the similarity coordinates
\NSFormula{NS-F-P007-015}{display}{3.2}
\begin{equation}
 \tau=q(1-\eta^2),\qquad z=q^D\eta,\qquad X=\frac{r^2}{2q},
 \qquad q>0,\qquad -1<\eta<1.
 \tag{3.2}\label{eq:3.2}
\end{equation}
For \NSi{NS-F-P007-016}{\tau>0}, removing \NSi{NS-F-P007-017}{\eta} from these equations leaves
\NSFormula{NS-F-P007-018}{display}{none}
\[
 q-z^2q^{2h}=\tau,\qquad
 \partial_q(q-z^2q^{2h})=1-2h\eta^2\geq 1-2h>0
 \quad\text{on }|\eta|<1.
\]
The derivative is positive, so these equations determine exactly one
\NSi{NS-F-P007-019}{q=q(z,\tau)}. They also give
\NSi{NS-F-P007-020}{q\asymp\tau+|z|^{1/D}}. Here
\NSi{NS-F-P007-021}{\asymp} means that fixed positive factors give bounds in both directions.
On any fixed compact set of similarity coordinates that stays away from
\NSi{NS-F-P007-022}{\eta=\pm1}, the more specific comparison
\NSi{NS-F-P007-023}{q\asymp\tau} holds. If \NSi{NS-F-P007-024}{X} stays bounded, letting
\NSi{NS-F-P007-025}{q\downarrow0} takes the coordinates toward the singular point. The endpoints
\NSi{NS-F-P007-026}{\eta=\pm1}, with \NSi{NS-F-P007-027}{q>0}, instead describe
\NSi{NS-F-P007-028}{t=1} away from that point, as Lemma~\ref{lem:4.1} shows.
\end{EESegment}

\begin{EESegment}{EE-TGT-P007-008}
The azimuthal and axial profiles \NSi{NS-F-P007-029}{E(X,\eta),U(X,\eta)} give the leading fields
through
\NSFormula{NS-F-P007-030}{display}{none}
\[
 u_\theta^{(0)}=q^{-A}E,\qquad u_z^{(0)}=q^{-A}U,
 \qquad
 ru_r^{(0)}=V_0,\qquad p^{(0)}=q^{-2A}\Pi.
\]
Incompressibility and regularity at the axis then fix \NSi{NS-F-P007-031}{V_0} through
\NSFormula{NS-F-P007-032}{display}{none}
\[
 \frac1r\partial_r(ru_r^{(0)})+\partial_z u_z^{(0)}=0,
 \qquad V_0(0,\eta)=0.
\]
\end{EESegment}
\NSPage{008}%
\begin{EESegment}{EE-TGT-P008-001}
The leading radial pressure balance, together with the normalization at radial infinity, gives
\NSFormula{NS-F-P008-001}{display}{none}
\[
 \partial_r p^{(0)}=\frac{(u_\theta^{(0)})^2}{r},
 \qquad
 \Pi(X,\eta)=-\int_X^\infty\frac{E(x,\eta)^2}{2x}\,dx.
\]
The construction makes \NSi{NS-F-P008-002}{E/\sqrt{2X}},
\NSi{NS-F-P008-003}{U}, and \NSi{NS-F-P008-004}{V_0/X} smooth at
\NSi{NS-F-P008-005}{X=0}. These are the combinations that have to be smooth for the Cartesian
field to be smooth at the axis whenever \NSi{NS-F-P008-006}{t<1}. In particular,
\NSi{NS-F-P008-007}{u_\theta^{(0)}=0} on the axis, while
\NSi{NS-F-P008-008}{E(X,\eta)>0} when \NSi{NS-F-P008-009}{X>0}.
\end{EESegment}

\begin{EESegment}{EE-TGT-P008-002}
Choose \NSi{NS-F-P008-010}{X_c>0} and \NSi{NS-F-P008-011}{0<\eta_c<1} so that the rectangle
\NSi{NS-F-P008-012}{[0,X_c]\times[-\eta_c,\eta_c]} stays inside the inner region from
Theorem~\ref{thm:4.6}. At time \NSi{NS-F-P008-013}{t=1-\tau}, define the core by
\NSFormula{NS-F-P008-014}{display}{none}
\[
 \mathcal C_\tau=\{(r\cos\theta,r\sin\theta,z):0\leq X(r,z,\tau)\leq X_c,
 \ |\eta(z,\tau)|\leq\eta_c\}.
\]
Throughout this core, \NSi{NS-F-P008-015}{q\asymp\tau}, so its radial and axial sizes satisfy
\NSFormula{NS-F-P008-016}{display}{none}
\[
 \ell_r\asymp\tau^{1/2},\qquad
 \ell_z\asymp\tau^{1/2-h},\qquad
 \frac{\ell_z}{\ell_r}\asymp\tau^{-h}.
\]
Both sizes go to zero, but their ratio goes to infinity.
\end{EESegment}

\begin{EESegment}{EE-TGT-P008-003}
The profiles give the following sizes for the three velocity components inside the core.
\NSFormula{NS-F-P008-017}{display}{none}
\[
 \begin{gathered}
  \|u_\theta^{(0)}\|_{L^\infty(\mathcal C_\tau)}\asymp\tau^{-1/2-h},
  \qquad
  \|u_z^{(0)}\|_{L^\infty(\mathcal C_\tau)}\asymp\tau^{-1/2-h},\\
  \|u_r^{(0)}\|_{L^\infty(\mathcal C_\tau)}=O(\tau^{-1/2}),
  \qquad
  \frac{\|u_r^{(0)}\|_{L^\infty(\mathcal C_\tau)}}
       {\|u_\theta^{(0)}\|_{L^\infty(\mathcal C_\tau)}}=O(\tau^h).
 \end{gathered}
\]
The lower bound for the axial component comes from the nonzero value on the axis in
Proposition~\ref{prop:B.2}. For any fixed \NSi{NS-F-P008-018}{0<X_*<X_c}, the circle given by
\NSi{NS-F-P008-019}{z=0,\ r=\sqrt{2X_*\tau}} has
\NSi{NS-F-P008-020}{q=\tau,\eta=0}, and on that circle
\NSFormula{NS-F-P008-021}{display}{none}
\[
 u_\theta^{(0)}=E(X_*,0)\tau^{-1/2-h}\longrightarrow+\infty.
\]
The core is a region of space that changes with time, while its boundary always has the same
similarity coordinates. A material region consists of the same fluid particles as the velocity
carries them along. Incompressibility keeps the volume of every such region fixed. Figure~\ref{fig:3}(a)
shows the radial regions.
\end{EESegment}

\begin{EESegment}{EE-TGT-P008-004}
The leading field continues outside the core. Outside a fixed interval of values of
\NSi{NS-F-P008-022}{X}, both \NSi{NS-F-P008-023}{U=V_0=0}, while
\NSFormula{NS-F-P008-024}{display}{none}
\[
 E(X,\eta)=c_\infty X^{-1/2-h}\bigl(1+O(X^{-1})\bigr)
 \qquad (X\to\infty),\qquad c_\infty>0,
\]
with one uniform estimate for every \NSi{NS-F-P008-025}{\eta}. The exact exterior formula
\eqref{eq:4.29} makes the exterior field a purely azimuthal solution of the radial heat equation
for the azimuthal velocity. The final spatial cutoff in Proposition~\ref{prop:10.1} leaves the
completed velocity supported inside one fixed compact region.
\end{EESegment}

\begin{EESegment}{EE-TGT-P008-005}
At viscosity one, the local radial and axial length scales
\NSi{NS-F-P008-026}{\ell_r=q^{1/2},\ \ell_z=q^D} give
\NSFormula{NS-F-P008-027}{display}{none}
\[
 \operatorname{Re}_\theta=q^{-A}\ell_r=q^{-h}\longrightarrow\infty,
 \qquad
 \operatorname{Re}_r=|u_r^{(0)}|\ell_r=O(1).
\]
The rotation rate is larger than the rate of radial diffusion. In the axisymmetric tangential
equations, radial diffusion and transport in both the radial and axial directions enter the leading
balance at the following rates.
\NSFormula{NS-F-P008-028}{display}{none}
\[
 \frac{|u_r^{(0)}|}{\ell_r}=O(q^{-1}),
 \qquad
 \frac{q^{-A}}{\ell_z}=\frac1{\ell_r^2}=q^{-1}.
\]
Axial diffusion is smaller than radial diffusion by the factor
\NSi{NS-F-P008-029}{\ell_r^2/\ell_z^2=q^{2h}}. This small factor is the expansion parameter that
orders the corrections to the background.
\end{EESegment}

\NSPage{009}%
\begin{figure}[htbp]
 \centering
 \NSFigurePlaceholder{NS-FIG-3}
 \begin{EESegment}{EE-TGT-P009-001}
 \caption{The concentrating geometry and the reference pulse. At one fixed time, panel (a) shows
 a radial--axial section. The pulse annulus lies between the two level surfaces
 \NSi{NS-F-P009-001}{X=X_a} and \NSi{NS-F-P009-002}{X=X_b}. The horizontal and vertical axes are
 each measured on their own similarity scale, and both physical length scales shrink. Panel (b)
 shows a reference amplitude \NSi{NS-F-P009-003}{P(v)}. It is normalized to equal one at the
 midpoint of a pulse interval of length \NSi{NS-F-P009-004}{L_s} in
 \NSi{NS-F-P009-005}{v}, and it first grows and then decays. The shaded parts contain the small
 tails made by the cutoff. The radii and pulse parameters in the picture are only representative.
 The proof chooses their actual values from the estimates in the sections that build the profiles
 and pulses.}
 \label{fig:3}
 \end{EESegment}
\end{figure}

\begin{EESegment}{EE-TGT-P009-002}
\subsection{Construction of the background stress.}\label{sec:3.2}
Next, choose the profiles \NSi{NS-F-P009-006}{E,U} so that the pulses can supply the remaining
momentum transport. The radial integrals below define the resulting
stress \NSi{NS-F-P009-007}{T}. This stress has to vanish outside the pulse annulus, and inside the
annulus it has to be a positive combination of the fluxes carried by the two wave families.
Theorem~\ref{thm:4.6}(ii) gives the required support. The admissible stress cone condition from
Lemma~\ref{lem:4.5} is the condition that makes the positive representation in
Proposition~\ref{prop:7.5} possible.
\end{EESegment}

\begin{EESegment}{EE-TGT-P009-003}
Let \NSi{NS-F-P009-008}{R_\theta^{(0)},R_z^{(0)}} be the tangential momentum residuals of
\NSi{NS-F-P009-009}{(u^{(0)},p^{(0)})}. These residuals include radial viscosity and leave
out axial viscosity. Construct physical stress components
\NSi{NS-F-P009-010}{T=(T_{r\theta},T_{rz})} that satisfy
\NSFormula{NS-F-P009-011}{display}{none}
\[
 R_\theta^{(0)}=-(\partial_r+2/r)T_{r\theta},
 \qquad
 R_z^{(0)}=-(\partial_r+1/r)T_{rz}.
\]
Regularity at the axis fixes the constants of integration, giving
\NSFormula{NS-F-P009-012}{display}{none}
\[
 T_{r\theta}(r)=-\frac1{r^2}\int_0^r s^2R_\theta^{(0)}(s)\,ds,
 \qquad
 T_{rz}(r)=-\frac1r\int_0^r sR_z^{(0)}(s)\,ds.
\]
\end{EESegment}

\begin{EESegment}{EE-TGT-P009-004}
Require \NSi{NS-F-P009-013}{T} to be nonzero exactly in the annulus
\NSi{NS-F-P009-014}{X_a<X<X_b}, which is the same region as
\NSi{NS-F-P009-015}{\sqrt{2qX_a}<r<\sqrt{2qX_b}}, with
\NSi{NS-F-P009-016}{0<X_a<X_b}. Make the residual zero both in the inner region and in the heat
exterior, and impose the two moment identities
\NSFormula{NS-F-P009-017}{display}{none}
\[
 \int_0^\infty r^2R_\theta^{(0)}\,dr=0,
 \qquad
 \int_0^\infty rR_z^{(0)}\,dr=0.
\]
These identities make the stress zero beyond the exterior radius as well as near the axis. The
prescribed exterior moment identities give exactly these two zero integrals, by
Lemma~\ref{lem:A.8}. Matching the profiles and the five radial integrals from \eqref{eq:4.15}
keeps the exterior fields unchanged, by Lemma~\ref{lem:4.4}.
\end{EESegment}

\NSPage{010}%
\begin{EESegment}{EE-TGT-P010-001}
At each point of the annulus, the admissible stress cone condition lets the pulse construction
choose two wave families whose covariance vectors
\NSi{NS-F-P010-001}{v_1,v_2\in\mathbb R^2} are linearly independent. Each vector has two entries.
They record the azimuthal and axial momentum flux produced there by one unit of squared amplitude
from that wave family in the local background. The two vectors can produce the required stress
with positive weights,
\NSFormula{NS-F-P010-002}{display}{none}
\[
 T=c_1v_1+c_2v_2,\qquad c_1,c_2>0\quad\text{where }T\neq0.
\]
Because both coefficients are positive, \NSi{NS-F-P010-003}{T} lies inside the cone generated by
\NSi{NS-F-P010-004}{v_1,v_2}. This cone is the set of all nonnegative combinations of the two
vectors, and being inside it means that this representation uses both directions with strictly
positive coefficients. Proposition~\ref{prop:7.5} builds the representation from the admissible
stress cone condition on the profiles. Figure~\ref{fig:4} shows the same positive representation.
\end{EESegment}

\begin{figure}[htbp]
 \centering
 \NSFigurePlaceholder{NS-FIG-4}
 \begin{EESegment}{EE-TGT-P010-002}
 \caption{Positive representation of the target stress. The shaded cone is the set of all
 nonnegative linear combinations of the covariance directions
 \NSi{NS-F-P010-005}{v_1,v_2}. The target
 \NSi{NS-F-P010-006}{T=c_1v_1+c_2v_2} is inside this cone because
 \NSi{NS-F-P010-007}{c_1,c_2>0}. The dashed lines show the two weighted vectors adding to
 the target stress. If the two directions change by a sufficiently small amount, both coefficients stay
 positive.}
 \label{fig:4}
 \end{EESegment}
\end{figure}

\begin{EESegment}{EE-TGT-P010-003}
These conditions tie the inner and exterior profiles together because both the pressure and the
radial velocity depend on cumulative radial integrals, the integrals collected from the axis out
to the current radius. For that reason, the pressure datum comes first. The equations are then
solved near the axis, and the integrals are matched before the cone condition is imposed.
\end{EESegment}
\begin{enumerate}[label=\arabic*.]
 \item \begin{EESegment}{EE-TGT-P010-004}
 \emph{Exterior and axis pressure.} Proposition~\ref{prop:A.4} builds reference profiles with a
 purely azimuthal tail. Their radial pressure integral \eqref{eq:4.31} fixes
 \NSi{NS-F-P010-008}{\Pi(0,\eta)}, the pressure value on the axis. Replacing this tail with the
 exact heat flow from Lemma~\ref{lem:A.6} changes the integral. The corrections in
 Proposition~\ref{prop:A.7} put back the same integral at smaller radii, so the pressure on the
 axis stays unchanged.
 \end{EESegment}

 \item \begin{EESegment}{EE-TGT-P010-005}
 \emph{Inner solution.} Using this pressure value and the specified axis data,
 Proposition~\ref{prop:B.2} constructs \NSi{NS-F-P010-009}{E/\sqrt{2X},U,\Pi} as functions that
 are analytic in \NSi{NS-F-P010-010}{X,\eta} near the axis, meaning that around each point,
 each function equals a convergent power series in those two variables, and satisfy \eqref{eq:4.13}. Their
 leading tangential residual and their stress are both zero on
 \NSi{NS-F-P010-011}{X\leq X_a}.
 \end{EESegment}

 \item \begin{EESegment}{EE-TGT-P010-006}
 \emph{Joining and moment matching.} Corollary~\ref{cor:B.10} joins the inner profiles to the
 prescribed outer profiles. The correction in Proposition~\ref{prop:B.8} makes the five
 cumulative radial integrals agree. These are the integrals that determine the pressure, radial
 velocity, and stress. When both the profiles and these integrals agree at the joining radius,
 Lemma~\ref{lem:4.4} gives back the prescribed exterior fields. The inner identity
 \eqref{eq:4.13} and the exterior conclusion of Lemma~\ref{lem:A.8} then give
 \NSi{NS-F-P010-012}{T=0} for \NSi{NS-F-P010-013}{X\leq X_a} and
 \NSi{NS-F-P010-014}{X\geq X_b}.
 \end{EESegment}

 \item \begin{EESegment}{EE-TGT-P010-007}
 \emph{Enforcing the cone condition.} On a compact subinterval of the annulus,
 Proposition~\ref{prop:C.2} adds a radial oscillation with phase
 \NSi{NS-F-P010-015}{N\log X}, where \NSi{NS-F-P010-016}{N} is one large fixed integer. The
 oscillation itself has amplitude \NSi{NS-F-P010-017}{O(N^{-1})}. Applying
 \NSi{NS-F-P010-018}{X\partial_X} removes that small factor from the derivative and produces a
 change of order one in the radial derivatives. This changes the leading tangential shear while
 changing the profile values and radial moments by only \NSi{NS-F-P010-019}{O(N^{-1})}. The
 shear follows a periodic curve that satisfies the admissible stress cone condition. A separate
 correction, supported in the same local region, restores all five radial integrals exactly. The
 inner solution and the heat exterior stay unchanged, and the resulting stress has positive
 covariance weights everywhere in the annulus, by Proposition~\ref{prop:C.3}.
 \end{EESegment}
\end{enumerate}

\NSPage{011}%
\begin{EESegment}{EE-TGT-P011-001}
The leading profiles are now fixed. Before the pulses are added, Section~\ref{sec:5} builds
corrections in the successive powers \NSi{NS-F-P011-001}{q^{2nh},\ n=1,2,\ldots}. At each order,
the coefficients already found provide the source terms in the equations for the next
coefficients. This recursion takes care of the radial momentum equation, axial viscosity, and all
the lower-order terms still left. Smooth cutoffs are then applied to the vector potentials and
pressures before the coefficients are added together. Each cutoff equals one near the singularity,
and the supports shrink from one correction to the next. Taking the curl after multiplication by
the cutoff keeps the field incompressible. The resulting smooth axisymmetric background
\NSi{NS-F-P011-002}{(u_B,p_B)} satisfies
\NSFormula{NS-F-P011-003}{display}{none}
\[
 \begin{aligned}
 \mathscr R(u_B,p_B)
   & =-(\partial_r+2/r)\mathcal T_{\mathrm{phys},\theta}e_\theta \\
   & \phantom{={}}-(\partial_r+1/r)\mathcal T_{\mathrm{phys},z}e_z+\mathcal E_B.
 \end{aligned}
\]
The stress \NSi{NS-F-P011-004}{\mathcal T_{\mathrm{phys}}} is supported in the annulus, and its
leading term is \NSi{NS-F-P011-005}{T}. Calling a residual \emph{flat} as
\NSi{NS-F-P011-006}{q\downarrow0} means that every Cartesian space--time derivative is
\NSi{NS-F-P011-007}{O(q^N)} for every \NSi{NS-F-P011-008}{N\geq0}, with one uniform estimate on
each bounded interval of \NSi{NS-F-P011-009}{X}. Proposition~\ref{prop:5.5} proves that the
remainder \NSi{NS-F-P011-010}{\mathcal E_B} is flat in this sense.
\end{EESegment}

\begin{EESegment}{EE-TGT-P011-002}
\subsection{Oscillations and momentum transport.}\label{sec:3.3}
At this point, the background \NSi{NS-F-P011-011}{(u_B,p_B)} has the required growth and the
leading annular stress \NSi{NS-F-P011-012}{T} that the two pulse families can produce. The next
step builds divergence-free pulses, using Lemma~\ref{lem:7.7}. When their quadratic flux is
averaged, its leading term supplies \NSi{NS-F-P011-013}{T}, by Propositions~\ref{prop:7.5}
and~\ref{prop:9.5}. The exact increment identity below keeps this cancellation separate from the
linear evolution of the pulses and from the interactions that later corrections must handle.
\end{EESegment}

\begin{EESegment}{EE-TGT-P011-003}
For a divergence-free velocity increment \NSi{NS-F-P011-014}{w} and a pressure increment
\NSi{NS-F-P011-015}{\pi}, the residual obeys the exact identity
\NSFormula{NS-F-P011-016}{display}{none}
\[
 \mathscr R(u_B+w,p_B+\pi)=\mathscr R(u_B,p_B)+\mathcal L_{u_B}(w,\pi)
   +\nabla\cdot(w\otimes w),
\]
where
\NSFormula{NS-F-P011-017}{display}{none}
\[
 \mathcal L_{u_B}(w,\pi):=\partial_t w+(u_B\cdot\nabla)w+(w\cdot\nabla)u_B
   -\Delta w+\nabla\pi.
\]
The linear term describes how the pulses evolve in the background. The quadratic term provides
the stress that cancels the leading part of
\NSi{NS-F-P011-018}{\mathscr R(u_B,p_B)}.
\end{EESegment}

\begin{EESegment}{EE-TGT-P011-004}
Write \NSi{NS-F-P011-019}{A_{\mathrm{wave}}} for the typical size of a pulse oscillation and
\NSi{NS-F-P011-020}{\ell_{\mathrm{wave}}} for its wavelength. At points in the pulse annulus
where the profile values are fixed and nonzero, set aside the logarithmic factors and the weights
that go to zero at the edge of the support. The remaining sizes are
\NSFormula{NS-F-P011-021}{display}{none}
\[
 \begin{gathered}
  A_{\mathrm{wave}}\asymp q^{-1/2-h/2},\qquad
  \ell_{\mathrm{wave}}\asymp q^{1/2+h/2},\\
  \frac{A_{\mathrm{wave}}}{q^{-1/2-h}}\asymp
  \frac{\ell_{\mathrm{wave}}}{q^{1/2}}\asymp q^{h/2}.
 \end{gathered}
\]
\end{EESegment}

\begin{EESegment}{EE-TGT-P011-005}
The square of the amplitude has the required stress scale.
\NSFormula{NS-F-P011-022}{display}{none}
\[
 A_{\mathrm{wave}}^2\asymp q^{-1-h}\asymp
 \frac{|u_\theta^{(0)}|}{q^{1/2}},
 \qquad
 \frac{A_{\mathrm{wave}}^2}{q^{1/2}}\asymp q^{-3/2-h}.
\]
The last quantity is the size of the averaged radial stress divergence. It matches the leading
time derivative, transport, and radial viscosity in the tangential momentum equations.
\end{EESegment}

\begin{EESegment}{EE-TGT-P011-006}
Freeze the amplitude and the phase gradient at one point. The leading wave then has the following
schematic local form.
\NSFormula{NS-F-P011-023}{display}{none}
\[
 w_{\mathrm{lead}}(x)=a\cos(\xi\cdot x+\varphi),\qquad a\cdot\xi=0.
\]
Here \NSi{NS-F-P011-024}{a\in\mathbb R^3} is the constant amplitude,
\NSi{NS-F-P011-025}{\xi\in\mathbb R^3} is the wavevector, and
\NSi{NS-F-P011-026}{\varphi\in\mathbb R} is the phase. The displayed orthogonality condition says
that the velocity oscillates in a direction perpendicular to its wavevector. Different localized
pulses have to stay separate.
\end{EESegment}
\NSPage{012}%
\begin{EESegment}{EE-TGT-P012-001}
At the same time, the construction still needs periodic averaging. It handles both needs by
introducing an extended field
\NSi{NS-F-P012-001}{\widetilde w(r,\theta,z,t,Y)} that depends on a separate auxiliary variable
\NSi{NS-F-P012-002}{Y\in\mathbb T^2}. The physical field comes from putting the fixed phase map
into that variable,
\NSFormula{NS-F-P012-003}{display}{none}
\[
 w(r,\theta,z,t)=\widetilde w(r,\theta,z,t,Y(r,t)),
\]
where \NSi{NS-F-P012-004}{Y(r,t)} is defined in \eqref{eq:6.3}. The covariance is averaged before
this substitution is made. The angular and auxiliary averages are
\NSFormula{NS-F-P012-005}{display}{none}
\[
 \langle F\rangle_\theta:=\frac1{2\pi}\int_0^{2\pi}F\,d\theta,
 \qquad
 \langle F\rangle_Y:=\int_{\mathbb T^2}F\,dY,
 \qquad
 \int_{\mathbb T^2}1\,dY=1.
\]
Doing the two averages one after the other is written as
\NSi{NS-F-P012-006}{\langle F\rangle:=\langle\langle F\rangle_\theta\rangle_Y}. Equations
\eqref{eq:3.7}--\eqref{eq:3.9} collect these conventions. All components are taken in the
cylindrical basis. If different labeled pulses overlap in their slowly varying space--time
supports, their supports in \NSi{NS-F-P012-007}{Y} are chosen to be disjoint. Their products are
then zero at every point, even after \NSi{NS-F-P012-008}{Y(r,t)} is substituted, by
Lemma~\ref{lem:6.1}. Harmonics and corrections with the same label still interact with one another.
\end{EESegment}

\begin{EESegment}{EE-TGT-P012-002}
In the actual construction, the amplitudes, phases, and cutoffs all vary. First evaluate the
localized vector potentials at \NSi{NS-F-P012-009}{Y(r,t)} to obtain
\NSi{NS-F-P012-010}{\mathcal A_{\mathrm{wave}}}. Then take the physical curl,
\NSFormula{NS-F-P012-011}{display}{none}
\[
 w=\nabla\times\mathcal A_{\mathrm{wave}},\qquad \nabla\cdot w=0.
\]
This curl makes the resulting velocity divergence-free. Every chain-rule term coming from
\NSi{NS-F-P012-012}{Y(r,t)} is included. Differentiating the amplitudes and cutoffs also produces
smaller velocity terms, and Lemma~\ref{lem:7.7} keeps all of them.
\end{EESegment}

\begin{EESegment}{EE-TGT-P012-003}
The chosen phases have nonzero integer angular frequencies. Because of that, each wave has angular
average zero, while the angular average of the square of its cosine is
\NSi{NS-F-P012-013}{1/2}. The two wave families produce the momentum-flux directions used above
to choose \NSi{NS-F-P012-014}{c_1,c_2>0}. Those coefficients then become the squared amplitude
weights. After the physical scaling and the corrections made by taking the curl, the covariance is
\NSFormula{NS-F-P012-015}{display}{none}
\[
 \begin{pmatrix}
  \langle\widetilde w_r\widetilde w_\theta\rangle\\
  \langle\widetilde w_r\widetilde w_z\rangle
 \end{pmatrix}
 =T+\text{higher-order terms}.
\]
\end{EESegment}

\begin{EESegment}{EE-TGT-P012-004}
The higher-order terms include the corrections made by the curl. They are smaller by positive
powers of \NSi{NS-F-P012-016}{q}, with derivative bounds proved in
Corollary~\ref{cor:7.8}. Taking the radial divergence of the leading covariance shown above
cancels the leading negative stress divergence of the background, by
Propositions~\ref{prop:7.5} and~\ref{prop:9.5}. The extended residual still contains oscillating
terms and angular means that depend on \NSi{NS-F-P012-017}{Y}. In the physical equation, these
terms are evaluated at \NSi{NS-F-P012-018}{Y(r,t)}.
\end{EESegment}

\begin{EESegment}{EE-TGT-P012-005}
The same quadratic products also control the energy passed from the background to the waves.
Take the homogeneous linearized model, where the linearized equation has zero right-hand side,
\NSFormula{NS-F-P012-019}{display}{none}
\[
 \mathcal L_{u_B}(w_{\mathrm{lin}},\pi_{\mathrm{lin}})=0,
 \qquad \nabla\cdot w_{\mathrm{lin}}=0.
\]
If the fields decay fast enough in space, integration by parts gives
\NSFormula{NS-F-P012-020}{display}{none}
\[
 \frac d{dt}\frac12\int|w_{\mathrm{lin}}|^2
 =-\sum_{i,j}\int(w_{\mathrm{lin}})_i(w_{\mathrm{lin}})_j\partial_j(u_B)_i
   -\int|\nabla w_{\mathrm{lin}}|^2.
\]
The first term moves energy between the background and the perturbation. The second term
dissipates energy. The chosen transverse velocity amplitudes draw energy from the shear and
provide the two momentum-flux directions whose positive span contains
\NSi{NS-F-P012-021}{T}. During each pulse, the shear increases the size of the radial component of
the wavevector. Viscous damping then grows until it is larger than the amplification, so the pulse
first grows and then decays. The time cutoff acts only in the exponentially small tails. The phase
and amplitude equations prove these facts in Lemmas~\ref{lem:7.1} and~\ref{lem:7.4}, and
Figure~\ref{fig:3}(b) shows the resulting envelope.
\end{EESegment}
